\section{Online link detection}
\label{app:linkdet}

Cross-document attention masks (\Cref{sec:method}) need to know, for each
generated or packed document, \emph{where} a link ends and \emph{what} it points
at. This appendix documents the online link detectors that supply that
information. The design goal is a detector that runs cheaply in the training hot
path and inside the generation loop, and that decouples link recognition from the
tokenizer so the mask machinery never has to know how text is encoded.

\subsection{The \texttt{LinkDetector} protocol}

A detector exposes up to three entry points. \texttt{detect\_links} scans a whole
one-dimensional token sequence and returns a list of \texttt{LinkInfo} records; it
runs at generation time over a rolling window (a link fires when its end position
coincides with the end of the window) and as the training fallback when no
per-document method exists. \texttt{detect\_links\_for\_doc} is the preferred
training path: it operates on a single document slice, with span-local positions
re-based by the caller, and because it receives the owning document's identifier it
can resolve \emph{relative} imports. It is not part of the base protocol --- the
caller probes for it with \texttt{hasattr} --- and is implemented only by the
detectors whose targets are relative (Python, TS/JS, Rust, Zig, Dart) and by the
composite detector. \texttt{index\_doc\_span} returns the \emph{lookup key} for a
document on the target side of an edge; the default is the raw identifier.
Detection matches a source-side target string against these keys by exact
dictionary lookup, so both sides must land in the same string space
(\Cref{app:linkdet:norm}). The traversal DAG additionally requires the target
document to start strictly before the link position (\Cref{app:traversal}).

The load-bearing design choice is that a \texttt{LinkInfo} carries a
\emph{decoded} target string and an \textbf{exclusive} end position
\texttt{link\_end\_pos} --- the token index just past the link's closing delimiter,
from which attention to the target is granted. The detector performs all decoding
and hands the mask creator plain strings plus integer offsets; because the offset
is an integer rather than a token identifier, the mask machinery has no tokenizer
dependency at all.

\subsection{Two detection strategies}

\paragraph{Token-id matching (markdown).}
Markdown links are detected \emph{without} a full decode. The detector searches for
the \texttt{](} bigram, which the GPT-2 BPE tokenizer emits as a single token
(id~$16151$), backtracks at most $100$ tokens for the opening \texttt{[} (ids $58$
or $685$), then grows a forward window (up to $50$ tokens), re-decoding the
accumulated chunk each step until a \texttt{)} appears. The target is the text
before the first \texttt{)} and \texttt{link\_end\_pos} is one past that token. This
is fast but tokenizer-specific (\Cref{app:linkdet:fail}).

\paragraph{Decode-once + blank + regex + remap (everything else).}
Every other detector decodes the span once, then blanks out comments and string
bodies by overwriting them with an equal-length run of spaces so that all remaining
character offsets are preserved. It then runs language-specific compiled regexes
over the blanked text and maps the matched character offsets back to token positions
via a cumulative per-token character-length array and \texttt{bisect\_left}. Python
takes a fast path that batch-decodes token bytes in one call (roughly $2\times$
faster on large files). This pipeline is a cheap approximation of a real parse:
offset-preserving blanking with per-language hand-written lexers handles nested
block comments, raw and byte strings, lifetimes-vs-char literals, and
template/multiline string bodies without a full grammar. Every runtime link
detector---for all languages and all three regimes below---uses this regex/token
pipeline; \emph{no} link detector calls a parser. Tree-sitter grammars
\citep{brunsfeld2018treesitter} are used elsewhere in the system: to extract the
\emph{shipped} graph offline, as an independent oracle that grades the runtime
detector for agreement, and---at benchmark-evaluation time only---to re-anchor the
scored span to the first use site of an imported symbol
(\Cref{app:linkdet:regimes}). Link \emph{detection} itself is grammar-free
throughout.

\subsection{Per-modality detectors}

\begin{itemize}
  \item \textbf{Markdown wikilinks:} token-id matching as above; target is the
        link title (spaces preserved).
  \item \textbf{Python:} three regexes (\texttt{import}, \texttt{from}-import, and a
        parenthesized multi-line form) with alias stripping, over
        comment/string/docstring-blanked text. It emits \emph{candidate file paths}
        rather than module names.
  \item \textbf{LaTeX \texttt{\textbackslash cite}:} a purely syntactic regex
        \texttt{\textbackslash cite[a-zA-Z]*(\lbrack...\rbrack)*\{...\}}; the bibkey is
        resolved to a title at extraction time, so detection returns the brace
        contents verbatim. Empty \texttt{\textbackslash cite\{\}} is skipped, and
        there is \emph{no} comment blanking.
  \item \textbf{Compiled-regex code detectors:} Go, Java, Kotlin, TS/JS, Zig, Dart,
        and Rust, all sharing the decode-blank-regex-remap pipeline. Their import
        syntax, key transforms, and blanking rules are summarized in
        \Cref{tab:linkdet}.
\end{itemize}

\subsection{Identifier normalization and key matching}
\label{app:linkdet:norm}

Matching is exact dictionary lookup, so a detector's target string and the graph's
\texttt{index\_doc\_span} key must be normalized into a common space. Some
modalities use \emph{identity} keys (the target string already equals the node
identifier); others transform the identifier. \Cref{tab:linkdet} lists the
per-language transform. A recurring subtlety is dual-mode resolution: for TS/JS the
flat scan produces specifier-space keys usable only by the detection grader, while
the per-document path produces resolved repository-relative keys that training
matches --- so the flat scan alone yields \emph{empty} masks for those languages.
Rust and Python similarly resolve \texttt{self}/\texttt{super} and relative imports
only in the per-document path.

\begin{table}[t]
  \centering
  \small
  \begin{tabular}{@{}llll@{}}
    \toprule
    Language & Import syntax detected & Key transform & Comment blanking \\
    \midrule
    Markdown   & \texttt{[title](...)} bigram & identity (title w/ spaces) & none \\
    arXiv/\texttt{\textbackslash cite} & \texttt{\textbackslash cite\{key\}} & identity (byte-identical title) & none \\
    Python     & \texttt{import}/\texttt{from} & bare-path $\to$ file candidates & \texttt{\#}, strings, docstrings \\
    Go         & single + grouped import & identity (full import path) & \texttt{//}, \texttt{/*\,*/}, strings/runes \\
    Java       & \texttt{import [static] a.b.C[.*];} & FQN dotification (drop \texttt{.java}) & none \\
    Kotlin     & \texttt{import a.b.C [as X]} & FQN (node id = symbol FQN) & \texttt{//}, \texttt{/*\,*/}, char, triple-quote \\
    Rust       & \texttt{use}/\texttt{mod} tree & strip \texttt{owner/repo@} $\to$ \texttt{crate::} path & comments, strings, char (not lifetimes) \\
    TS/JS      & \texttt{from}, dynamic \texttt{import()}, \texttt{require()} & extension-less path & comments + template bodies \\
    Zig        & \texttt{@import("...")} & keep \texttt{.zig} & \texttt{//} only, char, \texttt{\textbackslash\textbackslash} multiline \\
    Dart       & first \texttt{import}/\texttt{export} URI & keep \texttt{.dart} & anchored \texttt{(?m)(?:\^{}|;)} \\
    \bottomrule
  \end{tabular}
  \caption{Per-modality detector, import syntax, target$\leftrightarrow$key
  transform, and blanking. ``Identity'' keys require no transform; the others
  rewrite the detected string into the graph's node-id space. Rust node ids are
  \texttt{owner/repo@crate::a::b}, using \texttt{@} as the crate/path separator
  because module paths already contain \texttt{::}; the key is the \texttt{crate::}
  module path after the \texttt{@}. Java drops the repository prefix and
  \texttt{.java} and replaces \texttt{/} with \texttt{.}, but the source root is
  unknown (\Cref{app:linkdet:fail}).}
  \label{tab:linkdet}
\end{table}

Two target-side subtleties matter for correctness. Python's bare-path key is
expanded by \texttt{module\_path\_to\_file\_paths} into \texttt{foo/bar/baz.py} plus
\texttt{\_\_init\_\_.py} candidates, and where a name could be a submodule or a
symbol both candidates are emitted; resolution keeps whichever is a real node. Rust
performs the analogous submodule-vs-symbol expansion and use-tree flattening
(\texttt{use a::\{b, c::D\}} yields one key per leaf). Java \texttt{import static}
resolves to the enclosing type and \texttt{.*} wildcards resolve to nothing; Kotlin
drops wildcards; Go stores the verbatim import path with no expansion.

\subsection{The composite detector}

The merged multi-source model uses a composite detector over eleven members. Its
rule is to pick \emph{exactly one} sub-detector per document and never merge
results, which avoids cross-firing between languages. On the target side,
\texttt{index\_doc\_span} sniffs by identifier (\texttt{::}~$\to$ Rust; strip
repository prefix and map the extension; a hostname heuristic for Go), falling back
to identity for ambiguous title-like identifiers (wiki, arXiv). The per-document
path sniffs by identifier and then by content; the flat \texttt{detect\_links}
(generation, no identifier available) sniffs by content only. The content sniff
scores each candidate as a weighted sum of signature matches (\texttt{\textbackslash
cite\{}~$3.0$, \texttt{@import(}~$3.0$, \texttt{use::}~$2.0$, markdown~$1.0$, so a
stray \texttt{](} never outvotes a real code signature), distinguishes TS from JS by
marker tokens, breaks ties by a fixed priority (markdown last), and returns nothing
for prose. The safety argument is soft: a mis-sniff usually produces a target that
matches no real key and therefore no-ops. Mixed-syntax single documents, which would
need per-token routing, are explicitly deferred.

\subsection{Three regimes: training, benchmark evaluation, generation}
\label{app:linkdet:regimes}

Which detection path is exercised, and whether a parser is involved at all, differs
across the three settings in which links are resolved.

\begin{itemize}
  \item \textbf{Training.} The packer bypasses text detection entirely: the resolved
        graph edge is baked into the precomputed example, so grants come from the
        known edge list rather than from re-scanning tokens. The online regex detector
        is only a fallback for the live (non-precomputed) path. No parser runs.
  \item \textbf{Benchmark evaluation.} Grants are produced by the language's regex
        detector (the same token-space pipeline as training), but tree-sitter is used
        \emph{separately} to decide \emph{where} the completion is scored: the
        \texttt{use\_line} scope re-anchors scoring to the first line that references
        a symbol declared in a granted cross-file document, located with a per-language
        top-level-declaration query. This sidesteps import syntax and \texttt{import~*}
        (\Cref{app:results}). Detection is still regex; only the scoring \emph{scope}
        uses a grammar.
  \item \textbf{Generation.} There are no baked edges, so the regex detector (the
        composite, for a merged model) runs over the model's own output on a rolling
        window and fires when a link closes at the just-emitted token; a resolved link
        then triggers the link-following retrieval step. No parser runs.
\end{itemize}

\subsection{Limitations}
\label{app:linkdet:fail}

The detectors trade grammar-completeness for hot-path speed, and several of the
resulting approximations are worth stating plainly.
\begin{itemize}
  \item \textbf{Tokenizer coupling of the markdown path.} The \texttt{](} detector
        is bound to a specific GPT-2 token id; under a different tokenizer it detects
        nothing. Any \texttt{](} in prose or code fires (there is no comment
        blanking), and the $50$-token forward window truncates unusually long links.
  \item \textbf{arXiv title brittleness.} Matching requires the cited title to be
        byte-identical to the extracted title; whitespace, Unicode, or BPE drift
        breaks it silently, and \texttt{\textbackslash cite} inside a \texttt{\%}
        comment is still detected (no blanking).
  \item \textbf{Char-to-token remap.} The \texttt{bisect} remap is exact for ASCII
        source but may be off by one token on a rare multibyte token; the resulting
        \texttt{link\_end\_pos} error is at most one position.
  \item \textbf{Python resolution gaps.} Candidate file-path generation can
        over-produce; the repository-root${=}$path-root assumption fails for
        site-packages and \texttt{sys.path} manipulation; dynamic and
        \texttt{TYPE\_CHECKING} imports go undetected.
  \item \textbf{Java source-root ambiguity.} The FQN key cannot be verified against
        the builder's source root, and there is no comment blanking (detection relies
        on the \texttt{;} terminator and a line anchor).
  \item \textbf{Composite dispatch.} The no-op-on-mismatch safety is soft rather than
        guaranteed, and single-language-per-document dispatch does not handle a
        document that mixes link syntaxes.
  \item \textbf{Grant cap.} A per-example cap on cross-document grants (default $64$)
        silently drops the excess after sorting by link position.
\end{itemize}
